\section{和が一になる関数}\label{節:相補重み}

正弦と余弦の節では,反転した積分と加えると分子の和が分母に一致した.原点について対称な区間でも,同じことが起こる分数がある.\,$a>0$,\,$g$ を偶関数,\,$h$ を奇関数とすると,それは
\[
I=\int_{-a}^a\frac{g(x)}{1+\e^{h(x)}}\dd x
\]
の形をしている.指数関数を含む分母を取り除き,\,$\int_0^a g(x)\dd x$ に帰着させよう.

ここでも先に結果を述べ,その後で理由を確かめる.
\begin{equation}\label{式:偶関数と相補重み}
\int_{-a}^a\frac{g(x)}{1+\e^{h(x)}}\dd x
 =\int_0^a g(x)\dd x
\end{equation}
が成り立つ.指数関数の因子を一度も積分せずに,式から取り除けるのだ.

これを示そう.\,$w(x)=\frac{1}{1+\e^{h(x)}}$ と置く.\,$h$ は奇関数なので $h(-x)=-h(x)$ であり,
\[
\begin{aligned}
w(-x)&=\frac{1}{1+\e^{-h(x)}}\\
 &=\frac{\e^{h(x)}}{\e^{h(x)}+1}
\end{aligned}
\]
となる.分母は $w(x)$ と同じで,分子の和 $1+\e^{h(x)}$ が分母に一致するから,
\[
w(x)+w(-x)=1
\]
である.一方,\,$x=-t$ と置くと,\,$g$ が偶関数であることから
\[
I=\int_{-a}^a g(x)w(-x)\dd x
\]
となる.2式を加えると
\[
\begin{aligned}
2I&=\int_{-a}^a g(x)\{w(x)+w(-x)\}\dd x\\
 &=\int_{-a}^a g(x)\dd x\\
 &=2\int_0^a g(x)\dd x
\end{aligned}
\]
を得る.両辺を $2$ で割れば式\eqref{式:偶関数と相補重み}である.

\begin{bookpoint}
\textbf{POINT}\par
$x$ を $-x$ に変えた積分を作ってもとの積分と加え,\,$\frac{1}{1+\e^{h(x)}}+\frac{1}{1+\e^{-h(x)}}=1$ を使って指数関数の因子を消す.
\end{bookpoint}

\begin{example}{指数関数と対称性}
次の定積分を求めよ.
\[
I=\int_{-2}^2\frac{x^4+\cos x}{1+\e^{x^3+x}}\dd x.
\]
\end{example}

\paragraphbf{方針}
分子 $x^4+\cos x$ は偶関数,指数 $x^3+x$ は奇関数である.本節の式で $g(x)=x^4+\cos x$,\,$h(x)=x^3+x$,\,$a=2$ とした場合にあたる.

\paragraphbf{解答}
$I=\frac{32}{5}+\sin2$

\paragraphbf{解法}
$g(x)=x^4+\cos x$ は偶関数で,\,$h(x)=x^3+x$ は奇関数である.\,$x$ を $-x$ に変えると分子はそのまま,指数だけが負になる.したがって
\[
\frac{1}{1+\e^{h(x)}}+\frac{1}{1+\e^{-h(x)}}=1
\]
より
\[
2I=\int_{-2}^2(x^4+\cos x)\dd x
\]
となる.よって
\[
\begin{aligned}
I&=\int_0^2(x^4+\cos x)\dd x\\
 &=\left[\frac{x^5}{5}+\sin x\right]_0^2\\
 &=\frac{32}{5}+\sin2
\end{aligned}
\]
である.分子は区間全体で正である必要はない.式\eqref{式:偶関数と相補重み}に使われているのは,分子が偶関数であることだけである.

\paragraphbf{和が一になる関数の発想}
式\eqref{式:偶関数と相補重み}の出発点は,次の2つの式の比較である.
\[
\frac{1}{1+\e^x},\qquad
\frac{1}{1+\e^{-x}}=\frac{\e^x}{1+\e^x}
\]
2つの分母は等しく,分子の和 $1+\e^x$ が分母に一致するので,2つを足すと $1$ になる.\,$w(x)=\frac{1}{1+\e^x}$ は,対称な2点で値の和が $1$ になる因子なのである.\,$x$ を $-x$ に変えて元の式と加えると,この因子が掛かった部分は $w(x)+w(-x)=1$ に置き換わり,消えてしまう.必要なのは指数が $x$ であることではなく,反転すると指数の符号が変わることである.そのため,指数を奇関数 $h(x)$ に広げても同じ議論が使えた.

\paragraphbf{分子が奇関数の場合}
一方,分子が奇関数のときは,同じ結論にはならない.たとえば $a>0$ として
\[
K=\int_{-a}^a\frac{x}{1+\e^x}\dd x
\]
を式\eqref{式:左右の積分}で左右の組にすると,
\[
\begin{aligned}
K&=\int_0^a x\left(\frac{1}{1+\e^x}-\frac{1}{1+\e^{-x}}\right)\dd x\\
 &=\int_0^a\frac{x(1-\e^x)}{1+\e^x}\dd x
\end{aligned}
\]
となる.被積分関数は $0<x\le a$ で負なので,\,$K<0$ である.「分子の奇関数を対称区間で積分するから $0$」と考えると,分母が偶関数でないことを見落とす.積全体の対称性か,反転後の式との関係を確かめる必要がある.

\paragraphbf{中心を移す場合}
この考えは原点以外にも移せる.区間 $[a,b]$ の中央を $c=\frac{a+b}{2}$ とし,\,$g(a+b-x)=g(x)$,\,$w(x)+w(a+b-x)=1$ なら
\[
\int_a^b g(x)w(x)\dd x=\frac12\int_a^b g(x)\dd x
\]
である.たとえば $w(x)=\frac{1}{1+\e^{x-c}}$ はこの条件を満たす.\,$a+b-x-c=-(x-c)$ だからである.変数を $u=x-c$ に移せば,原点に関する場合と全く同じ証明になる.
